%% kapa: exam=1 round=23 year=2012 period=1 form=A subject=civil_law subject_ko=민법 seq=37 number=37
%% source: Q-Net 감정평가사 기출문제 (cst00309) artlSeq=5052578 file=raw/1008/5052578/2012년 제23회 감정평가사 제1차시험 기출문제(1교시 A형).pdf
%% conversion: kordoc 4.24.51 -> tools/gen1.py (auto); review=proofread-sonnet; flags=none
\begin{question}{37}
가등기담보권의 실행에 관한 설명으로 옳지 않은 것은? (다툼이 있으면 판례에 의함)
\begin{choices}
  \item 후순위권리자는 청산기간 내에 한하여 그 피담보채권의 변제기가 되기 전이라도 목적부동산의 경매를 청구할 수 있다.
  \item 채권자가 채무자에게 청산금을 통지하지 아니한 경우에도 청산금을 지급하고 등기를 마쳤다면 소유권을 취득할 수 있다.
  \item 채권자는 사적실행의 방법으로 청산기간이나 동시이행관계를 인정하지 아니하는 처분정산형의 담보권 실행을 할 수 없다.
  \item 목적부동산의 평가액이 채권액에 미달하여 청산금이 없다고 인정되는 때에는 채권자는 그 뜻을 채무자 등에게 통지하여야 한다.
  \item 압류등기 전에 이루어진 담보가등기권리가 매각에 의하여 소멸되면 경매법원에 채권신고를 한 경우에만 그 채권자는 매각대금을 배당받거나 변제금을 받을 수 있다.
\end{choices}
\end{question}
\answer{2}
